\chapter{The exact dyadic mixed-sign boundary}\label{ch:tree}
\section{Sectors and two different notions of separation}
For a packet put $c_i(v)=(v_i-i)/5$. At depth $k\ge1$, write $M=2^k$. A sector is the residue of $c(v)$ modulo $M$, identifying the residues obtained from $v$ and $Jv$. One may represent it by the lexicographically smaller of $c(v)\bmod M$ and $c(Jv)\bmod M$. All sector sums retain the original packet multiplicities.

A sector is mixed if it contains both positive and negative weights. Sign-purity is weaker than resolving a sector into one $J$-orbit. Several different orbits can have the same sign. A theorem about the first mixed sector therefore answers a different question from a theorem about the first collision of distinct orbits.

The sign has an elementary coding description. For $i<j$, the difference
\[
v_i-v_j=(i-j)+5(c_i-c_j)
\]
is positive exactly when $c_i>c_j$. There are ten pairs, so
\begin{equation}\label{eq:inversionsign}
\sgn W(v)=(-1)^{\#\{i<j:c_i>c_j\}}.
\end{equation}
Equal $c_i,c_j$ give a negative difference because $i-j<0$. Thus ties in the coding do not make a zero Vandermonde.

\section{The sharp threshold}
\begin{theorem}\label{thm:threshold}
Let $k\ge1$, $M=2^k$, and set
\begin{equation}\label{eq:nk}
N_k=\left\lceil\frac{100M^2+10M+2}{35}\right\rceil
=\frac{20M^2+2M+\gamma_k}{7},
\end{equation}
where
\[
\gamma_k=\begin{cases}
7&k\equiv1\pmod3,\\
1&k\equiv2\pmod3,\\
6&k\equiv0\pmod3.
\end{cases}
\]
Every depth-$k$ sector in every shell $n<N_k$ is sign-pure. The shell $n=N_k$ has a mixed sector.
\end{theorem}
This threshold ranges over all positive integer shell indices. It does not assert that $N_k$ is prime. For example, $N_3=186$ is composite.

\section{The geometric lower bound}
Take opposite-sign packets $v,w$ in one sector. Replace $w$ by $Jw$ if needed so that
\[
w-v=5Md,\qquad d\in A_4\setminus\{0\}.
\]
Set $a=5M/2$ and $m=(v+w)/2$. Since $M$ is even, both are integral. Equal norms give
\begin{equation}\label{eq:midpoint}
m\cdot d=0,\qquad10n=\|m\|^2+a^2\|d\|^2.
\end{equation}
The squared norm of a nonzero integral zero-sum vector is a positive even integer. It cannot be $2$ here. Such a vector is $\pm(e_i-e_j)$, and $m\cdot d=0$ would force $v_i-v_j=\pm5M$. This is divisible by $5$, contrary to the distinct coordinate residues.

If $\|d\|^2\ge6$, \eqref{eq:midpoint} gives
\[
n\ge\frac{15}{4}M^2>\frac{100M^2+10M+2}{35},\qquad M\ge2.
\]
The remaining case is $\|d\|^2=4$. Then $d$ has two coordinates $1$, two $-1$ and one $0$. Of the ten coordinate pairs, two have equal $d$ entries and eight have unequal entries. For a pair with $\delta=d_i-d_j$,
\[
(v_i-v_j)(w_i-w_j)=(m_i-m_j)^2-a^2\delta^2.
\]
The two equal-entry pairs have positive products. If all eight unequal-entry pairs had negative products, their product would also be positive, contradicting the opposite signs of $D(v)$ and $D(w)$. None of these pair products is zero, since both packets have distinct coordinates. Hence at least one unequal pair has a positive product, giving the integral inequality
\[
|m_i-m_j|\ge a|\delta|+1.
\]

Project $e_i-e_j$ orthogonally to $d^\perp$. Its squared norm is
\[
2-\frac{\delta^2}{4}.
\]
Because $m\perp d$, Cauchy's inequality gives
\[
\|m\|^2\ge\frac{(a|\delta|+1)^2}{2-\delta^2/4}.
\]
For $|\delta|=1$ this is $4(a+1)^2/7$. For $|\delta|=2$ it is $(2a+1)^2$, a larger quantity. Therefore
\[
10n\ge4a^2+\frac47(a+1)^2,
\]
which simplifies to the lower bound inside the ceiling in \eqref{eq:nk}. Taking the ceiling proves that no earlier shell can be mixed. Since powers of $2$ modulo $7$ cycle as $2,4,1$, direct division gives the displayed $\gamma_k$ values.

\section{Witnesses attaining the bound}
At $k=1$, take
\[
v=(-5,-4,2,-2,9),\qquad w=(5,6,-8,-2,-1).
\]
They belong to $\mathcal V_{13}$, have opposite weights and differ coordinatewise by multiples of $10$. Thus they lie in one depth-one sector.

For $k\ge2$, explicit affine families attain the lower bound. The following table lists $7v$ and $7w$, so that the displayed entries remain small integer affine expressions.
\begin{center}\small
\begin{tabular}{c>{\raggedright\arraybackslash}p{0.39\linewidth}>{\raggedright\arraybackslash}p{0.39\linewidth}}
\toprule
$k\bmod3$&$7v$&$7w$\\
\midrule
1, $k\ge4$&$(-20M+5,-10M-8,-10M-1,20M+16,20M-12)$&$(15M+5,-10M-8,25M-1,-15M+16,-15M-12)$\\[3pt]
2, $k\ge2$&$(-25M-5,10M+2,-15M+4,15M-4,15M+3)$&$(10M-5,10M+2,20M+4,-20M-4,-20M+3)$\\[3pt]
0, $k\ge3$&$(-20M-15,10M-3,10M+4,20M+1,-20M+13)$&$(15M-15,-25M-3,10M+4,-15M+1,15M+13)$\\
\bottomrule
\end{tabular}\end{center}
The appropriate powers of $2$ make each entry divisible by $7$. After division, each vector has the standard residues, sum zero and norm squared $10N_k$. The norm identities before division are
\[
\sum_i(7v_i)^2=\sum_i(7w_i)^2
=1400M^2+140M+70\gamma_k=490N_k.
\]
In the first two rows, $(w-v)/(5M)=(1,0,1,-1,-1)$. In the third it is $(1,-1,0,-1,1)$. Thus the packets share a sector.

Every coordinate difference is affine in $M$. Its sign is fixed on the indicated range. Checking the ten linear differences in each row proves that the two Vandermonde signs are opposite throughout that range. This is an all-depth argument using polynomial identities and linear inequalities, not an inference from a list of sample powers.

The first thresholds are
\[
13, 47, 186, 737, 2935, 11722, 46849, 187319.
\]
The verifier checks witnesses through $k=64$ as an implementation control. The affine identities and lower-bound proof establish the theorem for every depth.

\figplaceholder{mixed_threshold}{The first mixed shell $N_k$ at depths $1$ through $8$. These exact values belong to an all-depth geometric theorem. The index range is all positive integer shells, and the thresholds need not be prime.}

\section{A geometric theorem without a nonzero total}
Refining sectors eventually separates distinct $J$-orbits on a fixed finite shell. The sharp threshold describes how quickly opposite signs can share a sector. Neither assertion compares the total positive mass with the total negative mass. Geometric separation has not supplied an arithmetic inequality between them.

This is why the proposed sector-preserving finishing argument must be examined separately. Resolving signs can make cancellation visible. It does not remove cancellation if each resolved mass is then conserved.

\chapter{Why strict localization and automatic valuation fail}\label{ch:failures}
\section{Conserved sectors can preserve opposite signs}
Partition a shell into sectors $S_r$, with masses $t_r=\sum_{v\in S_r}W(v)$. Suppose every permitted correction is null inside every sector:
\[
\sum_{v\in S_r}W(v)Q_j(v)=0.
\]
Then $F=1+\sum a_jQ_j$ leaves each corrected sector mass equal to $t_r$. If two sector masses have opposite nonzero signs, all corrected packets cannot have one weak sign. Summing within those two sectors proves the obstruction immediately.

At $p=7$, depth-one sector data are
\begin{center}\small
\begin{tabular}{lrrr}
\toprule
Sector key&Packets&Mass&$\sum WQ_8$\\
\midrule
$(0,0,0,0,0)$&2&$-1848$&$-67060224000000$\\
$(1,0,0,1,0)$&2&$-17248$&$65383718400000$\\
$(1,1,0,0,0)$&2&2352&1676505600000\\
\bottomrule
\end{tabular}\end{center}
The three correction sums total zero, but no individual one is zero. Thus the known universal null $Q_8$ is not sector-null when restricted by a parity indicator. The sign certificate at $7$ works by transferring mass between sectors. A completion that preserves every one of these sectors would forbid the needed transfer.

This is a failure of the proposed localization principle, not a failure of the universal null identity. Multiplying $Q$ by an arbitrary sector indicator changes its theta series and may destroy universal nullness. A correct localized operation must be constructed by solving for combinations whose parent effects cancel, as in Chapter~\ref{ch:transport}.

\section{Aggregate parity identities have a larger scope}
The number of odd coordinates of a packet is even. Since squares modulo $4$ are $0$ or $1$,
\[
\#\{i:v_i\text{ odd}\}\equiv10n\equiv2n\pmod4.
\]
For odd $n$ there are two odd coordinates. For even $n$ there are zero or four. Let $A,B,C$ be the weighted generating series for these three aggregate classes, respectively.

The all-even class scales by $2$. The bijection $u_j=v_{2j\bmod5}/2$ has a coordinate permutation of sign $-1$, so
\[
A(q)=-2^{10}\Delta(q^4).
\]
Using the Hecke recurrence at $2$ gives
\begin{align}
B(q)&=-24\Delta(q^2)-2^{10}\Delta(q^4),\\
C(q)&=\Delta(q)+24\Delta(q^2)+2^{11}\Delta(q^4).
\end{align}
For a homogeneous universally null correction, scaling kills its aggregate $A$ series, and odd and even support then kill the corresponding aggregate $B$ and $C$ series. This does not imply nullity inside each of the six finer parity sectors making up class $C$. The explicit $p=7$ table is a counterexample to that stronger inference.

At sufficient depth, a sector is one $J$-orbit. For an even symmetric $Q$, its weighted correction sum on the orbit is its nonzero mass times $Q(v)$. Sector nullity then forces $Q(v)=0$ there. Imposing sector-nullity at every depth eventually removes every effective even correction. Refinement adds constraints rather than creating new available transfers.

\section{A unique least valuation is a sufficient criterion}
If one summand in an integer sum has strictly smaller $\ell$-adic valuation than all others, the sum has that least valuation and is nonzero. This follows by dividing by the least power of $\ell$ and observing one unit plus terms divisible by $\ell$.

Terminal packet resolution does not guarantee this hypothesis. Every packet is $1$ modulo $5$, and each $J$-orbit has size $1$ or $2$. Every orbit mass is therefore a $5$-adic unit. At an index prime $p>20$, Proposition~\ref{prop:indexunit} makes every orbit mass a $p$-adic unit as well. When several orbits occur, there is no unique least valuation at either of these primes.

At $p=31$, depth $2$ already resolves all sixteen $J$-orbits. Eight masses are positive and eight negative. Four have the minimum $2$-adic valuation $3$. At depth $1$, two sector masses tie at that minimum. The unique-minimum test therefore fails at every positive depth in this example, although
\[
\tau(31)=-52843168,\qquad v_2(\tau(31))=5.
\]
Further residue information about tied sums may be useful. The counterexample excludes the proposed automatic criterion, rather than every possible valuation argument.

\section{A prime-index interpolation that cannot be integral}
An earlier proposed shortcut treated
\[
F(u)=\frac{\tau(2^{14}u-1)}{2^{14}}
\]
at prime arguments as the restriction of a single integral $2$-adic power series. Any convergent integral restricted power series is $1$-Lipschitz: subtracting its values at $x,y$ factors by $x-y$ with an integral convergent quotient.

Take the prime indices
\[
(u_1,p_1)=(15,245759),\qquad(u_2,p_2)=(47,770047).
\]
The exact values are
\begin{align*}
\tau(245759)&=77900595298037436891536424960,\\
\tau(770047)&=-379872617393156753656047295381504.
\end{align*}
Their residues modulo $2^{19}$ are $393216$ and $147456$. Dividing by $2^{14}$ gives
\[
F(15)\equiv24\pmod{32},\qquad F(47)\equiv9\pmod{32}.
\]
Thus $v_2(F(47)-F(15))=0$, while $v_2(47-15)=5$. This violates the necessary Lipschitz condition and refutes the stated integral interpolation model.

Different valuations alone would not disprove analyticity. The obstruction is this divided-difference failure. The example does not exclude models with denominators, finer residue classes or additional variables. It also does not concern the legitimate auxiliary polynomial $G_n(U)$ at a fixed shell. The two constructions use different variables and different hypotheses.

\section{Exact lifting equalities also need proof}
Another proposed shortcut was
\[
v_\ell(\tau(p))=v_\ell(p+1)
\]
for candidate primes $p\ne\ell$ with $p\equiv-1\pmod\ell$, at $\ell=2,3,5$. All three equalities fail. Exhaustive search through $150000$ gives the first failures within these respective candidate domains:
\begin{center}
\begin{tabular}{rrrr}
\toprule
$\ell$&$p$&$v_\ell(\tau(p))$&$v_\ell(p+1)$\\
\midrule
2&81919&20&14\\
3&8747&8&7\\
5&499&4&3\\
\bottomrule
\end{tabular}
\end{center}
At $p=1249$, the $5$-adic discrepancy goes the other way: the valuations are $3$ and $4$. The decisive integers were independently checked through the Eisenstein convolution. These are counterexamples to the auxiliary valuation equalities, never to Lehmer. They show that higher cancellation cannot be read off from $p+1$ by the stated formula.
\figplaceholder{lifting_failures}{The first failures of the three proposed valuation equalities in their stated candidate domains, found by exhaustive prime search through $150000$. All decisive values were independently checked. The equalities fail, while Lehmer is not refuted.}
